% Original English questions and concise English solutions. Compile using XeLaTeX.
\documentclass[UTF8,10pt,letterpaper,fontset=fandol]{ctexart}
\usepackage[margin=1in]{geometry}
\usepackage{amsmath,amssymb}
\usepackage{enumitem,needspace}
\usepackage[unicode,hidelinks]{hyperref}
\pagestyle{plain}
\raggedbottom
\setlength{\parindent}{0pt}
\setlength{\parskip}{1pt}
\setlength{\emergencystretch}{2em}
\setcounter{secnumdepth}{0}
\setlist[enumerate]{label=(\arabic*),leftmargin=2em,itemsep=0.25em,topsep=0.2em}
\ctexset{section={format=\Large\bfseries,beforeskip=1.2ex,afterskip=0.8ex}}
\newcommand{\C}{\mathbb C}
\newcommand{\Solution}{\par\Needspace{4\baselineskip}\textbf{Solution}\par}
\begin{document}
\begin{center}
{\LARGE Homework 5}\\[0.6em]
October 21, 2026\\[0.25em]
Name: YuanHao Tang\quad Student ID: 3120265621
\end{center}

\section{Problem 1}
Let
\[
B(\lambda)=\begin{pmatrix}\lambda^3+\lambda&\lambda^2\\\lambda^2-3\lambda&2\lambda\end{pmatrix},\qquad
C(\lambda)=\begin{pmatrix}
1-\lambda&\lambda^2&\lambda\\
2\lambda-1&\lambda(1-\lambda)&-2\lambda\\
1+\lambda^2&\lambda^2&-\lambda^2
\end{pmatrix},
\]
and
\[
A(\lambda)=\begin{pmatrix}B(\lambda)&0\\0&C(\lambda)\end{pmatrix}.
\]
Find out
\begin{enumerate}
\item the Smith canonical form of $B(\lambda)$;
\item the elementary divisors of $A(\lambda)$.
\end{enumerate}
\Solution
Work over $\C[\lambda]$ and let $D_k$ denote the $k$th determinantal divisor.
\begin{enumerate}
\item Every entry of $B$ is divisible by $\lambda$, and $2\lambda$ is an entry, so $D_1(B)=\lambda$. Also
\[
\det B=\lambda^2(\lambda+1)(\lambda+2),
\qquad\boxed{S_B=\operatorname{diag}\bigl(\lambda,\lambda(\lambda+1)(\lambda+2)\bigr)}.
\]
\item For $C$, direct calculation gives $D_1=1$, $D_2=\lambda$ and $D_3=\lambda^2(\lambda+1)$, hence
\[
S_C=\operatorname{diag}\bigl(1,\lambda,\lambda(\lambda+1)\bigr).
\]
Take the multiset union of the elementary divisors of both blocks. All elementary divisors of $A$ are
\[
\boxed{\lambda,\ \lambda,\ \lambda,\ \lambda,\ \lambda+1,\ \lambda+1,\ \lambda+2}.
\]
\end{enumerate}

\section{Problem 2}
Let $D:\C[\lambda]\to\C[\lambda]$ be defined by
\[
D\!\left(\sum_{i=0}^n\alpha_i\lambda^i\right)=\sum_{i=1}^n i\alpha_i\lambda^{i-1}.
\]
Show that $D$ is a linear transformation on $\C[\lambda]$, and determine $\ker[D]$.
\Solution
Differentiation maps polynomials to polynomials and satisfies
$D(uf+vg)=uD(f)+vD(g)$ for polynomials $f,g$ and complex scalars $u,v$, so $D$ is linear.
Moreover, $D(f)=0$ iff $i\alpha_i=0$ for every $i\ge1$, which means $f$ is constant. Thus
\[
\boxed{\ker D=\C\cdot1=\operatorname{span}\{1\}}.
\]

\newpage
\section{Problem 3}
Let $\mathcal A\in L(\C_n[\lambda])$ be defined, for a nonzero constant $a\in\C$, as
\[
\mathcal A[f(\lambda)]=f(\lambda+a),\qquad\forall f(\lambda)\in\C_n[\lambda].
\]
Find out
\begin{enumerate}
\item the matrix representation of $\mathcal A$ with respect to the basis $1,\lambda,\lambda^2,\ldots,\lambda^n$;
\item the eigenvalues and the associated eigenvectors of $\mathcal A$.
\end{enumerate}
\Solution
\begin{enumerate}
\item By $(\lambda+a)^j=\sum_{i=0}^j\binom ji a^{j-i}\lambda^i$, the matrix $M=[\mathcal A]$, with rows and columns indexed from $0$ to $n$, is
\[
\boxed{M_{ij}=\begin{cases}\binom ji a^{j-i},&0\le i\le j\le n,\\0,&i>j.\end{cases}}
\]
\item $M$ is upper triangular with all diagonal entries one, so its only eigenvalue is $1$, of algebraic multiplicity $n+1$.
If $f$ has degree $m\ge1$ and leading coefficient $c_m$, then $f(\lambda+a)-f(\lambda)$ has nonzero leading term $mc_ma\lambda^{m-1}$.
Thus only constant polynomials satisfy $f(\lambda+a)=f(\lambda)$, and
\[
\boxed{E_1(\mathcal A)=\operatorname{span}\{1\},\qquad f(\lambda)=c\ (c\in\C\setminus\{0\})}.
\]
The nonzero constant polynomials are precisely all associated eigenvectors.
\end{enumerate}

\section{Problem 4}
Let $V=\C_n[s]$ and $D:\C_n[s]\to\C_n[s]$ be defined by
\[
\begin{aligned}
D[f(s)]&=na_ns^{n-1}+(n-1)a_{n-1}s^{n-2}+\cdots+2a_2s+a_1,\\
f(s)&=a_ns^n+a_{n-1}s^{n-1}+\cdots+a_2s^2+a_1s+a_0.
\end{aligned}
\]
Find out
\begin{enumerate}
\item the elementary divisors of $D$;
\item the eigenvalues and the associated eigenvectors of $D$.
\end{enumerate}
\Solution
\begin{enumerate}
\item In the basis $(1,s,\ldots,s^n)$, $N=[D]$ has nonzero entries only at $N_{j-1,j}=j$ for $1\le j\le n$.
Use $\mu$ as the characteristic-matrix indeterminate. The superdiagonal of $\mu I-N$ is $-1,-2,\ldots,-n$.
For each $1\le k\le n$ there is a nonzero constant minor $(-1)^k k!$, so $D_k=1$; also $\det(\mu I-N)=\mu^{n+1}$. Hence
\[
S_{\mu I-N}=\operatorname{diag}\bigl(\underbrace{1,\ldots,1}_{n},\mu^{n+1}\bigr),
\qquad\boxed{\text{the sole elementary divisor is }\mu^{n+1}}.
\]
\item $N$ is strictly upper triangular, so its only eigenvalue is $0$, of algebraic multiplicity $n+1$. Since $D(f)=f'=0$ exactly for constant polynomials,
\[
\boxed{E_0(D)=\operatorname{span}\{1\},\qquad f(s)=c\ (c\in\C\setminus\{0\})}.
\]
All associated eigenvectors are the nonzero constant polynomials.
\end{enumerate}
\end{document}
